\documentclass[11pt]{article}

\usepackage[margin=1in]{geometry}
\usepackage{amsmath,amssymb,amsthm,mathtools}
\usepackage{booktabs}
\usepackage{enumitem}
\usepackage{graphicx}
\usepackage{hyperref}
\usepackage{xcolor}
\usepackage{microtype}

\hypersetup{colorlinks=true,linkcolor=blue!55!black,urlcolor=blue!55!black,citecolor=blue!55!black}
\setlist{nosep}
\newtheorem{theorem}{Theorem}
\newtheorem{lemma}{Lemma}
\newcommand{\R}{\mathbb{R}}
\newcommand{\cE}{\mathcal{E}}
\newcommand{\cF}{\mathcal{F}}
\newcommand{\supp}{\operatorname{supp}}
\newcommand{\norm}[1]{\left\lVert #1\right\rVert}

\title{Smooth Kernels Suffice for Two-Layer\\Universal Operator Approximation}
\author{Candidate full solution for arXiv:1910.01545}
\date{August 11, 2026}

\begin{document}
\maketitle

\begin{center}
\fbox{\parbox{0.88\linewidth}{\textbf{Status: full solution, likely valid, pending human review.}
The theorem below affirmatively answers the existence part of the first future-work
question in Section~6 of Guss--Salakhutdinov, and removes the suggested extra
regularity assumption on the target operator.}}
\end{center}

\section{Source problem and scope}

The source is William H. Guss and Ruslan Salakhutdinov,
\emph{On Universal Approximation by Neural Networks with Uniform Guarantees on
Approximation of Infinite Dimensional Maps}, arXiv:1910.01545v1 (2019)~\cite{GS}.
On page~12, Section~6, the authors ask:

\begin{quote}
``Given some regularity conditions on a nonlinear operator, are the weight
functions of the operator layer based neural networks that approximate it
smooth or regular?''
\end{quote}

Figure~\ref{fig:source} reproduces the source context.  We answer the
existence question in the affirmative.  In fact, continuity of the nonlinear
operator and compactness of the input family already suffice; no additional
regularity of the operator is needed.

\begin{figure}[ht]
  \centering
  \includegraphics[width=0.98\linewidth]{figures/open_problem_crop.png}
  \caption{The first future-work question in Section~6, page~12 of
  arXiv:1910.01545v1.}
  \label{fig:source}
\end{figure}

The conclusion is qualitative and existential.  It does not provide bounds
on derivatives, sample complexity, or the outcome of a training algorithm.

\section{Set-up and exact statement}

Let $K\subset\R^d$ and $K'\subset\R^{d'}$ be the compact domains of the
source paper, equipped with Lebesgue measure.  Let $C(K)$ and $C(K')$ carry
the supremum norm.  We say that a function on a compact Euclidean set is
\emph{ambient smooth} if it is the restriction of a $C^\infty$ function on
the surrounding Euclidean space.  Our construction gives compactly supported
ambient extensions.

For a latent domain $I=[0,1]$, smooth kernels and biases define the two-layer
network
\begin{equation}\label{eq:operator-net}
 \mathcal{N}(f)(y)
 =\int_I W_2(v,y)\,
 g\!\left(\int_K f(u)W_1(u,v)\,du+B_1(v)\right)dv+B_2(y).
\end{equation}
This is precisely a composition of two affine operator layers with a
pointwise activation between them.

\begin{theorem}[smooth two-operator-layer universality]\label{thm:main}
Let $K\subset\R^d$ and $K'\subset\R^{d'}$ be compact domains, let
$E\subset C(K)$ be compact, and let $F:C(K)\to C(K')$ be continuous.  If
$g:\R\to\R$ is continuous and nonpolynomial, then for every $\varepsilon>0$
there exist
\[
 W_1\in C_c^\infty(\R^d\times\R),\qquad
 W_2\in C_c^\infty(\R\times\R^{d'}),
\]
and $B_1\in C_c^\infty(\R)$, $B_2\in C_c^\infty(\R^{d'})$ such that the
network \eqref{eq:operator-net}, using their restrictions to
$K\times I$, $I\times K'$, $I$, and $K'$, satisfies
\[
 \sup_{f\in E}\norm{\mathcal{N}(f)-F(f)}_\infty<\varepsilon.
\]
Consequently, two-layer operator networks whose two weight kernels and two
biases are jointly ambient $C^\infty$ are dense, on every compact input
family, in $C(C(K),C(K'))$.
\end{theorem}

\section{Proof intuition}

The key is that the source paper has already reduced the infinite-dimensional
problem to finitely many hidden features.  Its Theorem~2 gives
\[
 N_0(f)(y)=\sum_{i=1}^N a_i(y)
 g\!\left(\int_K f(u)c_i(u)\,du+b_i\right)+a_0(y),
\]
where $c_i\in L^1(K)$ and $a_i,a_0\in C(K')$.  On the compact input set $E$,
the input norms are uniformly bounded.  Thus small $L^1$ changes in $c_i$
make all preactivations uniformly close.  The activation is uniformly
continuous on the single compact interval containing both the original and
perturbed preactivations.  Smooth density in $L^1$ and Stone--Weierstrass
therefore smooth all coefficients without losing uniform accuracy.

The resulting expression is still a finite hidden sum.  To express it
\emph{exactly} with two integral operator layers, assign one disjoint smooth
bump to each hidden unit in the latent interval.  A companion cutoff is one
on that bump's support, so the latent integral reads each hidden feature
without cross-talk.  The integral of each normalized bump is one, hence the
integral is exactly the finite sum.

\section{Proof}

\begin{proof}[Proof of Theorem~\ref{thm:main}]
Apply Theorem~2 of the source paper~\cite{GS} with tolerance
$\varepsilon/2$.  There are $N\geq1$, functions
$c_i\in L^1(K)$, $a_i\in C(K')$, constants $b_i\in\R$, and
$a_0\in C(K')$ such that
\begin{equation}\label{eq:source-approx}
 \sup_{f\in E}\norm{N_0(f)-F(f)}_\infty<\frac{\varepsilon}{2},
\end{equation}
where
\[
 N_0(f)(y)=\sum_{i=1}^N a_i(y)g(z_i(f))+a_0(y),
 \qquad
 z_i(f)=\int_K f(u)c_i(u)\,du+b_i.
\]

\paragraph{Step 1: smooth the finite-hidden network.}
Compactness of $E$ gives
\[
 M:=\max_{f\in E}\norm{f}_\infty<\infty.
\]
Put
\[
 A:=\sum_{i=1}^N\norm{a_i}_\infty,
 \qquad
 R:=1+\max_{1\leq i\leq N}\bigl(M\norm{c_i}_{L^1}+|b_i|\bigr),
 \qquad
 G:=\max_{|t|\leq R}|g(t)|.
\]
By uniform continuity of $g$ on $[-R,R]$, choose $0<\eta<1$ so that
\begin{equation}\label{eq:uc}
 |g(s)-g(t)|<\frac{\varepsilon}{8(1+A)}
 \quad\text{if }s,t\in[-R,R]\text{ and }|s-t|<\eta.
\end{equation}

Extend each $c_i$ by zero outside $K$.  Density of
$C_c^\infty(\R^d)$ in $L^1(\R^d)$ gives $p_i\in C_c^\infty(\R^d)$ with
\begin{equation}\label{eq:input-smooth}
 \norm{p_i-c_i}_{L^1(K)}<\frac{\eta}{M+1}.
\end{equation}
For every $f\in E$,
\[
 \left|\int_K f(u)(p_i-c_i)(u)\,du\right|
 \leq M\norm{p_i-c_i}_{L^1(K)}<\eta.
\]
The original preactivations have absolute value at most $R-1$; hence the new
ones
\[
 z'_i(f):=\int_K f(u)p_i(u)\,du+b_i
\]
belong to $[-R,R]$.  Therefore \eqref{eq:uc} applies uniformly in $f$.

By the real Stone--Weierstrass theorem, polynomial restrictions are uniformly
dense in $C(K')$.  Multiplying the approximating polynomials by a smooth
cutoff equal to one on a neighborhood of $K'$ shows that restrictions of
$C_c^\infty(\R^{d'})$ are uniformly dense in $C(K')$.  Choose
$q_i,q_0\in C_c^\infty(\R^{d'})$ such that
\begin{equation}\label{eq:output-smooth}
 \norm{q_i-a_i}_{C(K')}<\frac{\varepsilon}{8N(1+G)},
 \qquad
 \norm{q_0-a_0}_{C(K')}<\frac{\varepsilon}{8}.
\end{equation}
Define the smoothed finite-hidden network
\[
 N_s(f)(y):=\sum_{i=1}^Nq_i(y)g(z'_i(f))+q_0(y).
\]
Adding and subtracting $a_i(y)g(z'_i(f))$, and using
\eqref{eq:uc}--\eqref{eq:output-smooth}, yields
\begin{align*}
 \norm{N_s(f)-N_0(f)}_\infty
 &\leq \sum_{i=1}^N\norm{q_i-a_i}_\infty G
  +\sum_{i=1}^N\norm{a_i}_\infty
      \frac{\varepsilon}{8(1+A)}
  +\norm{q_0-a_0}_\infty \\
 &<\frac{\varepsilon}{8}+\frac{\varepsilon}{8}
   +\frac{\varepsilon}{8}
 =\frac{3\varepsilon}{8}.
\end{align*}
Together with \eqref{eq:source-approx}, this gives
\begin{equation}\label{eq:smoothed-error}
 \sup_{f\in E}\norm{N_s(f)-F(f)}_\infty
 <\frac{7\varepsilon}{8}<\varepsilon.
\end{equation}

\paragraph{Step 2: encode the finite sum exactly by smooth operator layers.}
Choose pairwise disjoint open intervals $J_1,\dots,J_N$ compactly contained
in $(0,1)$.  For each $i$, choose a nonnegative
$\zeta_i\in C_c^\infty(J_i)$ with
$\int_0^1\zeta_i(v)\,dv=1$.  Choose
$\chi_i\in C_c^\infty(J_i)$ that equals one on $\supp\zeta_i$.  Set
\begin{align*}
 W_1(u,v)&:=\sum_{i=1}^N\chi_i(v)p_i(u),
 &B_1(v)&:=\sum_{i=1}^N\chi_i(v)b_i,\\
 W_2(v,y)&:=\sum_{i=1}^N\zeta_i(v)q_i(y),
 &B_2(y)&:=q_0(y).
\end{align*}
These functions have the asserted compactly supported smooth ambient
extensions.  If $v\in\supp\zeta_i$, disjointness gives
$\chi_i(v)=1$ and $\chi_j(v)=0$ for $j\neq i$.  Consequently,
\begin{align*}
 &\int_0^1 W_2(v,y)
 g\!\left(\int_K f(u)W_1(u,v)\,du+B_1(v)\right)dv+B_2(y)\\
 &\quad=\sum_{i=1}^Nq_i(y)
    \int_0^1\zeta_i(v)g\!\left(\int_Kf(u)p_i(u)\,du+b_i\right)dv+q_0(y)\\
 &\quad=\sum_{i=1}^Nq_i(y)g(z'_i(f))+q_0(y)
 =N_s(f)(y).
\end{align*}
The operator-layer network is therefore exactly $N_s$, and
\eqref{eq:smoothed-error} proves the claim.
\end{proof}

\section{Verification notes}

\begin{enumerate}
\item \textbf{Topology check.}  The only use of compactness of $E$ is the
uniform input bound $M$.  Continuity of $F$ and nonpolynomiality of $g$ enter
through the source paper's Theorem~2.
\item \textbf{Norm check.}  $L^1$ approximation of $c_i$ is the correct norm
because inputs are bounded in $L^\infty$; H\"older's inequality gives the
uniform preactivation bound.
\item \textbf{Activation check.}  No global Lipschitz assumption on $g$ is
used.  Both old and new preactivations lie in $[-R,R]$, where continuity
implies uniform continuity and boundedness.
\item \textbf{Joint smoothness check.}  Each kernel is a finite sum of
products of ambient smooth functions in separate variables, hence is jointly
$C^\infty$.  All supports are compact.
\item \textbf{Architecture check.}  The disjoint supports eliminate all
cross-terms, and normalization $\int\zeta_i=1$ makes the latent integral
exactly the finite hidden sum.
\item \textbf{Mechanical check.}  The script
\texttt{code/verify\_error\_budget.py} verifies that the allocated errors sum
to $7\varepsilon/8$ and checks the diagonal overlap algebra of the bump
encoding using exact rational arithmetic.
\end{enumerate}

\section{Limitations and novelty check}

The theorem answers whether smooth weights \emph{can be chosen} in uniformly
approximating operator-layer networks.  It does not show that every
representation is smooth, that training discovers such weights, or that an
exact representation of $F$ exists.  It is qualitative: no derivative bounds,
rates, network-size estimates, or sample bounds are claimed.  The other two
future directions in the source conclusion (variational optimization and
sub-exponential sampling bounds) are untouched.

A bounded novelty search was performed on August~11, 2026.  The run's cheap
registry, solution, attempt, and proof-gap indexes were searched for the arXiv
identifier and the core phrases.  The official arXiv API was queried for
``operator layer'' with ``smooth weight'', ``neural operator'' with ``smooth
kernel'' and ``universal approximation'', and the exact source title.  Bounded
arXiv web searches used close variants.  They found later neural-operator
universality results, notably Kovachki et al.~\cite{Kovachki} and a KNO theorem
with continuous positive-definite kernels~\cite{KNO}, but no theorem located
in those bounds states the present ambient-smooth, two-layer,
$C(K)\to C(K')$, uniform-on-compact-family conclusion.  This is a bounded
search, not an exhaustive literature review; novelty confidence remains
moderate.

\section{Human-review recommendation}

The argument is short, uses only standard density results, and preserves the
source architecture exactly.  Review should focus on (i) matching the source's
``compact domains'' convention for Lebesgue integration, (ii) confirming that
the future-work question is intended existentially rather than as a request
for quantitative regularity inherited from $F$, and (iii) checking for an
unindexed later statement of the same smoothing observation.  Mathematically,
the packet is recommended for full-solution review.

\begin{thebibliography}{9}
\bibitem{GS}
W. H. Guss and R. Salakhutdinov,
\emph{On Universal Approximation by Neural Networks with Uniform Guarantees on
Approximation of Infinite Dimensional Maps}, arXiv:1910.01545v1, 2019.

\bibitem{Kovachki}
N. Kovachki, Z. Li, B. Liu, K. Azizzadenesheli, K. Bhattacharya, A. Stuart,
and A. Anandkumar,
\emph{Neural Operator: Learning Maps Between Function Spaces},
arXiv:2108.08481.

\bibitem{KNO}
M. Lowery, J. Turnage, Z. Morrow, J. D. Jakeman, A. Narayan, S. Zhe, and
V. Shankar,
\emph{Kernel Neural Operators (KNOs) for Scalable, Memory-efficient,
Geometrically-flexible Operator Learning}, arXiv:2407.00809.

\bibitem{Rudin}
W. Rudin, \emph{Real and Complex Analysis}, 3rd ed., McGraw--Hill, 1987.
Standard references for $L^1$ approximation and Stone--Weierstrass.
\end{thebibliography}

\end{document}

